%% 2i2d-palais-sports-treillis — exercice guidé : ferme en treillis de 56 m de
%% portée posée sur deux poteaux béton. Un appui glissant (type A) laisse la
%% dilatation se produire, un appui fixe (type B) bloque la translation.
\documentclass[tikz,border=4pt]{standalone}
\usepackage{liaisons}
\begin{document}
\begin{tikzpicture}[x=1cm,y=1cm,
  barre/.style={line width=0.7pt, solideB, line cap=round},
  membrure/.style={line width=1.1pt, solideB, line cap=round},
  cote/.style={{Stealth[length=4pt]}-{Stealth[length=4pt]}, line width=0.6pt, solideD},
  attache/.style={line width=0.4pt, solideD!80},
  vent/.style={-{Stealth[length=5pt]}, line width=1pt, solideD},
  etiq/.style={font=\small, inner sep=2pt},
  petit/.style={font=\scriptsize, inner sep=1.5pt}]

\newcommand\batihoriz[3]{%
  \draw[line width=1pt, solideE] (#1,#3) -- (#2,#3);
  \begin{scope}
    \clip (#1,#3-0.24) rectangle (#2,#3);
    \foreach \hx in {-16,-15.76,...,16} { \draw[line width=0.5pt, solideE!75] (\hx,#3) -- (\hx-0.28,#3-0.28); }
  \end{scope}}

%% --- ordonnée de la membrure supérieure (deux pentes faibles, dissymétriques) ---
\newcommand\haut[1]{ifthenelse(#1<4.6, 1.75+0.2162*(#1-0.9), 2.55-0.2*(#1-4.6)) }

%% --- sol et poteaux béton --------------------------------------------------------
\batihoriz{0.2}{7.8}{0}
\hachures{(0.675,0) rectangle (1.125,1.2)}
\hachures{(6.875,0) rectangle (7.325,1.2)}
\node[petit, anchor=north] at (0.9,-0.28) {poteau 1};
\node[petit, anchor=north] at (7.1,-0.28) {poteau 2};

%% --- la ferme en treillis ----------------------------------------------------------
\draw[membrure] (0.9,1.2) -- (7.1,1.2);                       % membrure inférieure
\draw[membrure] (0.9,1.75) -- (4.6,2.55) -- (7.1,2.05);       % membrure supérieure
\foreach \x in {1.8,2.7,3.6,4.6,5.5,6.4} {                    % montants
  \draw[barre] (\x,1.2) -- (\x,{\haut{\x}}); }
\draw[barre] (0.9,1.2) -- (0.9,1.75)  (7.1,1.2) -- (7.1,2.05);
\foreach \a/\b in {0.9/1.8, 2.7/1.8, 2.7/3.6, 4.6/3.6, 4.6/5.5, 6.4/5.5, 6.4/7.1} {
  \draw[barre] (\a,1.2) -- (\b,{\haut{\b}}); }                % diagonales en zigzag

%% --- couverture et panneaux ----------------------------------------------------------
\draw[line width=1pt, solideF] (0.9,1.9) -- (4.6,2.7) -- (7.1,2.2);
\node[petit, anchor=south, text=solideF] at (3.2,2.72) {couverture + panneaux photovoltaïques};

%% --- barre 1 : membrure supérieure juste à droite du sommet, en traction ----------------
\draw[line width=1.3pt, solideA] (4.6,2.55) -- (5.5,2.37);
\draw[-{Stealth[length=4pt]}, line width=0.8pt, solideA] (5.05,2.46) -- (5.85,2.30);
\draw[-{Stealth[length=4pt]}, line width=0.8pt, solideA] (5.05,2.46) -- (4.25,2.62);
\node[petit, anchor=north, text=solideA] at (5.35,2.22) {barre 1 — $N = 1\,750$ kN};

%% --- appuis ----------------------------------------------------------------------------
\draw[line width=0.9pt, draw=solideE, fill=solideE!14, line join=round]
  (0.9,1.17) -- (0.66,0.86) -- (1.14,0.86) -- cycle;
\draw[line width=0.7pt, draw=solideE, fill=white] (0.78,0.75) circle (0.09);
\draw[line width=0.7pt, draw=solideE, fill=white] (1.02,0.75) circle (0.09);
\node[petit, anchor=east, text=solideE, align=right] at (0.55,0.62) {appui glissant\\(type A)};
\draw[line width=0.9pt, draw=solideE, fill=solideE!14, line join=round]
  (7.1,1.17) -- (6.86,0.86) -- (7.34,0.86) -- cycle;
\node[petit, anchor=west, text=solideE, align=left] at (7.45,0.72) {appui fixe\\(type B)};

%% --- le vent des deux côtés ---------------------------------------------------------------
\draw[vent] (0.05,3.05) -- (0.75,3.05);
\node[petit, anchor=south, text=solideD] at (0.4,3.08) {vent d'ouest};
\draw[vent] (7.95,3.05) -- (7.25,3.05);
\node[petit, anchor=south, text=solideD] at (7.6,3.08) {vent d'est};

%% --- la dilatation -------------------------------------------------------------------------
\draw[{Stealth[length=4pt]}-{Stealth[length=4pt]}, line width=0.7pt, solideD,
      dash pattern=on 3pt off 2pt] (3.2,0.95) -- (4.8,0.95);
\node[petit, anchor=north, text=solideD] at (4.0,0.92) {dilatation $\approx$ 27 mm};

%% --- portée ------------------------------------------------------------------------------------
\draw[attache] (0.9,-0.3) -- (0.9,-0.7)  (7.1,-0.3) -- (7.1,-0.7);
\draw[cote] (0.9,-0.55) -- (7.1,-0.55);
\node[petit, anchor=south, text=solideD] at (4.0,-0.52) {56 m};
\end{tikzpicture}
\end{document}
